\VerdierPhysicalPageAnchor{0115}%
\begingroup
\small
\setlength{\abovedisplayskip}{4pt plus 2pt minus 2pt}%
\setlength{\belowdisplayskip}{4pt plus 2pt minus 2pt}%

\VerdierSemanticAnchor{ch2-1-2-3-proof-f1-f2}
Supposons, pour fixer les idées, que $f_1\colon X\longrightarrow X'$ et
$f_2\colon Y\longrightarrow Y'$ soient des isomorphismes. Pour tout objet
$W$ de $\mathcal D$, on a un diagramme commutatif :

\VerdierSemanticAnchor{ch2-1-2-3-five-lemma-diagram}
\begin{equation*}
  \resizebox{0.98\linewidth}{!}{%
    $\begin{tikzcd}[ampersand replacement=\&,column sep=2.15em,row sep=3em]
      \cdots \arrow[r]
        \& \operatorname{Hom}_{\mathcal D}(W,X) \arrow[r] \arrow[d,"{\operatorname{Hom}_{\mathcal D}(W,f_1)}"']
        \& \operatorname{Hom}_{\mathcal D}(W,Y) \arrow[r] \arrow[d,"{\operatorname{Hom}_{\mathcal D}(W,f_2)}"']
        \& \operatorname{Hom}_{\mathcal D}(W,Z) \arrow[r] \arrow[d,"{\operatorname{Hom}_{\mathcal D}(W,f_3)}"']
        \& \operatorname{Hom}_{\mathcal D}(W,X[1]) \arrow[r] \arrow[d,"{\operatorname{Hom}_{\mathcal D}(W,f_1[1])}"']
        \& \cdots \\
      \cdots \arrow[r]
        \& \operatorname{Hom}_{\mathcal D}(W,X') \arrow[r]
        \& \operatorname{Hom}_{\mathcal D}(W,Y') \arrow[r]
        \& \operatorname{Hom}_{\mathcal D}(W,Z') \arrow[r]
        \& \operatorname{Hom}_{\mathcal D}(W,X'[1]) \arrow[r]
        \& \cdots
    \end{tikzcd}$%
  }\,.
\end{equation*}

\VerdierSemanticAnchor{ch2-1-2-3-five-lemma}
où les lignes sont exactes. On en déduit, par le lemme des cinq, que
$\operatorname{Hom}_{\mathcal D}(W,f_3)$ est un isomorphisme pour tout $W$,
et par suite que $f_3$ est un isomorphisme. Les deux autres assertions du
corollaire se déduisent de celle-ci en utilisant l'axiome TRII.

\VerdierSemanticAnchor{ch2-1-2-4}
\medskip
\noindent\textbf{Corollaire 1.2.4.}\quad
Soient
$X\xrightarrow{u}Y\xrightarrow{v}Z\xrightarrow{w}X[1]$ et
$X\xrightarrow{u}Y\xrightarrow{v'}Z'\xrightarrow{w'}X[1]$
deux triangles distingués. Il existe un isomorphisme
$f\colon Z\xrightarrow{\sim}Z'$ tel que le diagramme ci-après soit
commutatif :

\VerdierSemanticAnchor{ch2-1-2-4-1}
\begin{equation}
  \begin{tikzcd}[column sep=2.8em,row sep=1.7em]
    && Z \arrow[dr,"w"] \arrow[dd,"f"',"\sim"] & \\
    X \arrow[r,"u"] & Y \arrow[ur,"v"] \arrow[dr,"v'"'] && X[1] \\
    && Z' \arrow[ur,"w'"'] &
  \end{tikzcd}
  \tag{1.2.4.1}
\end{equation}

\VerdierSemanticAnchor{ch2-1-2-4-proof}
En effet, d'après l'axiome TRIII, il existe un morphisme
$f\colon Z\longrightarrow Z'$ tel que le diagramme (1.2.4.1) soit
commutatif. D'après (1.2.3), ce morphisme est un isomorphisme.

\VerdierSemanticAnchor{ch2-1-2-4-essential-uniqueness}
Le corollaire (1.2.4) montre donc que les triangles distingués contenant un
morphisme donné sont essentiellement uniques : ils sont déterminés à
isomorphisme près. Mais il convient de noter que cet isomorphisme n'est pas
unique. Nous préciserons ce point en (1.2.12).

\VerdierSemanticAnchor{ch2-1-2-5}
\medskip
\noindent\textbf{Corollaire 1.2.5.}\quad
La somme directe de deux triangles distingués est un triangle distingué.

\VerdierSemanticAnchor{ch2-1-2-5-proof-start}
Soient :

\VerdierSemanticAnchor{ch2-1-2-5-input-triangles}
\[
  \begin{gathered}
    X\xrightarrow{u}Y\xrightarrow{v}Z\xrightarrow{w}X[1],\\
    X'\xrightarrow{u'}Y'\xrightarrow{v'}Z'\xrightarrow{w'}X'[1]
  \end{gathered}
\]
\endgroup
